\section{Plano de análise e discussão}

A análise prevista reconstruirá trajetórias de indicadores por evento de divulgação e categoria, verificando continuidades, seleções, reconfigurações e divergências entre as arenas. A Tabela~\ref{tab:parametros-recorte} apresenta parâmetros do recorte, e não frequências observadas ou resultados da pesquisa. A comparação distinguirá ano-base do fenômeno, ano de publicação e edição ou revisão do produto institucional. Diferenças numéricas não serão automaticamente corrigidas: será necessário verificar se correspondem a classificação, denominador, revisão administrativa ou outra diferença documentável antes de estabelecer condições de comparação.

A discussão examinará em que medida os vínculos encontrados autorizam falar em circulação ou mobilização de indicadores, preservando os limites de cada evidência. A literatura sobre atenção política justifica observar relações entre indicadores e documentos estatais, sem antecipar o sentido dos resultados no caso brasileiro \cite{kristensen2023indicadores}. Nenhum teste estatístico, distribuição de frequências ou associação empírica é apresentado neste esqueleto. O produto técnico esperado é um Protocolo de Rastreabilidade, Comparabilidade e Comunicação Responsável de Indicadores de Violência Letal, a ser fundamentado na investigação posterior.

\begin{table}[htb]
\caption{Parâmetros do recorte proposto, sem resultados empíricos}
\label{tab:parametros-recorte}
\centering
\begin{tabular}{p{0.30\textwidth}p{0.60\textwidth}}
\toprule
\textbf{Parâmetro} & \textbf{Definição no projeto} \\
\midrule
Anos-base & 2016 a 2024 \\
Comunicação & Institucional e jornalística, em arenas distintas \\
Veículos & G1, Folha de S.Paulo e O Estado de S.Paulo \\
Mobilização estatal & Legislativo e Executivo federal \\
Inferência & Relações documentais, sem causalidade presumida \\
\bottomrule
\end{tabular}
\fonte{Elaboração do autor para o exercício (2026), a partir das delimitações do projeto.}
\end{table}
